\subsection{An actual imperfect family: function fields over
  \texorpdfstring{$\overline{\mathbf F}_p$}{the algebraic closure of a finite field}}
\label{C:function-fields-subsection}

Fix a prime $p$ and a coherent computable presentation
$C=\overline{\mathbf F}_p$ as the union of the finite fields
$\mathbf F_{p^{s!}}$. Embeddings between successive fields can
be found by finite search. Elements have finite-stage names,
and arithmetic and equality are computed in a common stage.
Keeping the first previously unseen name gives a computable
enumeration without repetition. Put $k=C(t)$, with rational
functions represented by coprime numerator and monic denominator.
All field extensions in this subsection are supplied as finite
towers of monic irreducible polynomial presentations. Their
iterated power bases provide exact arithmetic and enumeration.
The algorithms below prove the required factorization access;
it is not an additional oracle assumption.

\begin{lemma}[Proved: effective Frobenius roots in every finite tower]
\label{C:function-field-pth-roots}
For every explicitly presented finite extension $K/k$, membership
in $K^p$ is decidable, and the unique $p$-th root of an element
of $K^p$ is computable. The same holds for iterated $p$-th roots.
\end{lemma}
\begin{proof}
Write $k^p=C(t^p)$. The inverse of Frobenius on $C$ is computable
in any finite field containing the input. Hence the isomorphism
$k\longrightarrow k^p$, $a\mapsto a^p$, and its inverse are
effective: take $p$-th roots of the constant coefficients and
replace the variable $t^p$ by $t$.

The elements $1,t,\ldots,t^{p-1}$ form a basis of $k$ over $k^p$.
Their coordinates are computable. Indeed, if $r(t)=f(t)/g(t)$,
then
\[
 r(t)=\frac{f(t)g(t)^{p-1}}{g(t)^p};
\]
the denominator belongs to $C[t^p]$. Grouping the monomials of
the numerator by their exponent modulo $p$ writes $r$ in the
stated basis, with coefficients in $C(t^p)$.

Let $e_1,\ldots,e_n$ be the known $k$-basis of $K$ obtained from
its tower. Then
\[
 \{t^a e_i:0\leq a<p,\ 1\leq i\leq n\}
\]
is a computable $k^p$-basis of $K$. Moreover
\[
             K^p=\sum_{i=1}^n k^p e_i^p.               \tag{C.3.K1}
\]
This follows by taking $p$-th powers in the expansion of an
arbitrary element in the $e_i$. The $e_i^p$ are linearly
independent over $k^p$: a relation with coefficients $a_i^p$
is the $p$-th power of $\sum_i a_i e_i$, and Frobenius is
injective on a field.

For $z\in K$, compute $z$ and all $e_i^p$ in the displayed
$np$-element basis and solve the finite linear system
$z=\sum_i u_i e_i^p$ over $k^p$. Inconsistency proves
$z\notin K^p$. Otherwise
$\sum_i u_i^{1/p}e_i$ is its unique $p$-th root.
This is a finite computation, with no assumption that $K/k$
is separable. Repeating it gives the asserted iterated procedure.
\end{proof}

\begin{proposition}[Proved: a finite constant-field bound for base factorization]
\label{C:function-field-constant-bound}
Univariate polynomials over $k=C(t)$ admit a terminating
factorization algorithm. More explicitly, let a nonconstant
polynomial in $k[X]$ have primitive denominator-cleared form
$F(t,X)\in\mathbf F_q[t,X]$, of total degree $D\geq1$.
Every normalized absolutely irreducible factor of $F$ is defined
over $\mathbf F_{q^L}$, where
$L=\operatorname{lcm}(1,\ldots,D)$.
Thus a finite search over this one finite field suffices.
\end{proposition}
\begin{proof}
All input coefficients lie in one finite constant field
$\mathbf F_q$. Clearing denominators and removing the greatest
common divisor of the coefficients as polynomials in $t$
are effective Euclidean computations in $\mathbf F_q[t]$.
Fix a monomial order in $\mathbf F_q[t,X]$. Normalize each
irreducible factor in $C[t,X]$ by requiring its leading
coefficient to be $1$.

The $q$-power Frobenius permutes these normalized factors,
because it fixes $F$ and preserves the normalization.
There are at most $D$ distinct factors: every factor has
positive total degree, and total degrees add in a product.
Each Frobenius orbit therefore has length $r\leq D$.
The coefficients of a factor in an orbit of length $r$
are fixed by the $q^r$-power Frobenius, and so belong to
$\mathbf F_{q^r}$. Since every such $r$ divides $L$,
all factors lie in $\mathbf F_{q^L}[t,X]$.
This proves the bound, including repeated factors.

Construct a finite constant field containing $\mathbf F_{q^L}$
inside the chosen presentation of $C$. Enumerate its finitely
many normalized polynomials of total degree at most $D$.
Exact polynomial division tests which divide $F$.
Recursively divide by nonconstant proper factors until none
remain; every step lowers total degree. All tests and searches
are finite. The resulting factors are absolutely irreducible,
because the proved bound already includes every absolute factor.
Factors depending only on $t$ are units in $C(t)[X]$.
For the other factors, Gauss's lemma identifies irreducibility
in $C[t,X]$ with irreducibility over $C(t)$ after primitive
normalization. Make them monic in $X$ and recover multiplicities
by exact division. This factors the original polynomial.
The degree-zero cases and its nonzero scalar are handled directly.
\end{proof}

\begin{theorem}[Proved: ordinary factorization through arbitrary finite towers]
\label{C:function-field-factorization}
There is a terminating algorithm for univariate factorization
over every explicitly presented finite tower over $C(t)$.
It applies to inseparable tower steps and to inseparable
input polynomials. Consequently every such field has the
finite-tower factorization access required in
Theorem~\ref{C:artinian_effective_imperfect_C1}.
\end{theorem}
\begin{proof}
We give the tower induction, with the inseparable operations
specified explicitly. The base is
Proposition~\ref{C:function-field-constant-bound}, and every
field occurring is infinite, countable and effectively
enumerable. Lemma~\ref{C:function-field-pth-roots} supplies
Frobenius-root tests at every step, independently of this
factorization induction.

\emph{Reduction to separable input.}
Suppose first that one can factor separable polynomials over a
field $E$ in the tower. This suffices to factor every polynomial
over $E$, as follows. Work with monic nonconstant polynomials
and use induction on degree.
If $f'\ne0$ and $\gcd(f,f')\ne1$, factor this proper gcd and
the quotient recursively; both degrees are smaller.
If $\gcd(f,f')=1$, the polynomial is separable and the supposed
procedure applies.
If $f'=0$, write $f(X)=h(X^p)$ and factor the polynomial $h$,
whose degree is smaller, recursively.

For each monic irreducible factor $a(Y)$ of $h$, the polynomial
$a(X^p)$ is either irreducible or the $p$-th power of a monic
irreducible polynomial. To verify this assertion, let $\beta$
be a root of $a$, let $\gamma^p=\beta$, and put
$d=[E(\beta):E]$. The extension
$E(\gamma)/E(\beta)$ has degree $1$ or $p$, and all roots of
$a(X^p)$ have the same unique $p$-th-root construction from
roots of $a$. In the degree-$p$ case the minimal polynomial
of $\gamma$ already has degree $pd$, so equals $a(X^p)$.
In the degree-$1$ case it has degree $d$; equivalently, the
unique irreducible factor has multiplicity $p$.
For the latter uniqueness, any two roots of $a$ are conjugate
over $E$ after passage to an algebraic closure, and the
unique $p$-th roots are conjugate as well.
This argument remains valid when $a$ is inseparable.

The second case is decided and its factor found by asking
whether every coefficient of $a(X^p)$ has a $p$-th root in $E$.
When all do, take those roots and divide all exponents by $p$;
the resulting polynomial is the irreducible factor.
When one does not, $a(X^p)$ is irreducible.
The root tests terminate by
Lemma~\ref{C:function-field-pth-roots}.
The multiplicities from $h$ and this step are multiplied.
This proves the reduction, without confusing a polynomial
having no repeated irreducible factors with a separable
polynomial over an imperfect field.

\emph{A separable simple extension.}
Assume factorization over $E$ is already available, and let
$K=E(\alpha)$ be given by a monic separable irreducible
polynomial. To factor a monic separable $f\in K[X]$,
enumerate $c\in E$ and compute
\[
 N_c(X)=N_{K/E}\bigl(f(X-c\alpha)\bigr)\in E[X]
\]
until $\gcd(N_c,N_c')=1$.
The norm is a determinant on the known $E$-basis of $K$,
so this is effective.
This search terminates. In an algebraic closure of $E$,
its roots are $\beta+c\,\sigma(\alpha)$, where $\sigma$
ranges over the $E$-embeddings of $K$ and $\beta$ ranges
over the roots of $\sigma(f)$. Within a fixed embedding
these roots are distinct because $f$ is separable.
For different embeddings the values $\sigma(\alpha)$
are different, so any pair can collide for at most one
value of $c$. There are only finitely many pairs.
The infinite field $E$ therefore contains a suitable $c$,
which its enumeration eventually finds.

Factor $N_c$ by the already available procedure over $E$.
For each of its monic irreducible factors $q$, compute
\[
        \gcd_{K[X]}\bigl(f(X-c\alpha),q(X)\bigr).       \tag{C.3.K2}
\]
The nonconstant gcds are exactly the monic irreducible
factors of $f(X-c\alpha)$, with no repetitions.
Here is the norm argument. The norm of a monic irreducible
$K$-polynomial is a positive power of one monic irreducible
$E$-polynomial: every conjugate of one root has the same
minimal polynomial over $E$, and the product defining the
norm is invariant under all $E$-automorphisms of a normal
closure. Alternatively, apply the minimal polynomial to
a root and all its conjugates.
If the norm had two relatively prime nonconstant factors
over $E$, divisibility of one by the original irreducible
polynomial would force the same divisibility by every
conjugate, contradicting the other factor.
Since $N_c$ is separable, this positive power is one, and
two distinct factors of $f(X-c\alpha)$ cannot have the
same norm factor. Formula (C.3.K2) therefore recovers them
one at a time. Replace $X$ by $X+c\alpha$ to finish.
The preceding reduction handles all remaining inputs over $K$.

This is the norm-and-gcd construction of
Trager~\cite[Section 2, Theorems 2.1--2.3 and Corollary 2.5,
pp.\ 219--221]{C-TragerFunctionFields1976}.
That source states a characteristic-zero hypothesis.
The argument just given proves precisely the separable-input,
separable-extension version in the present characteristic;
the inseparable steps are not inferred from that hypothesis.

\emph{An inseparable simple extension.}
Now let $K=E(\alpha)$ have monic irreducible polynomial
$g(T)=h(T^{p^e})$, where $e\geq1$ and $h'\ne0$.
The polynomial $h$ is irreducible, for a factorization
would give one of $g$. Put $\beta=\alpha^{p^e}$ and
$E'=E(\beta)$. Then $E'/E$ is a separable simple extension
and $K/E'$ is purely inseparable of exponent at most $p^e$.
The previous step provides factorization over $E'$.
Its embedding in $K$ is explicit, and coordinate
conversion is a finite linear computation on power bases.
If $h$ has degree one, this simply means $E'=E$.

For monic $f\in K[X]$, the polynomial $f(X)^{p^e}$ has
coefficients in $E'$: every element of $K$ has its
$p^e$-th power in $E'$. Compute this polynomial and
factor it over $E'$.
Each monic irreducible factor $q\in E'[X]$ has exactly
one monic irreducible factor over $K$, with multiplicity
a power of $p$ dividing $p^e$. Indeed, if $u\in K[X]$
is a monic divisor of $q$, then $u^{p^e}\in E'[X]$
divides $q^{p^e}$. Irreducibility of $q$ gives
$u^{p^e}=q^b$ for an integer $b$.
Two different irreducible divisors would have relatively
prime powers, so this identity proves uniqueness.
Writing $q=v^a$ for its unique irreducible factor gives
$v^{p^e}=q^b=v^{ab}$, and thus $a\mid p^e$.

Starting with $q$, repeatedly take a polynomial $p$-th
root whenever all exponents are divisible by $p$ and
every coefficient has a $p$-th root in $K$.
The required membership tests and roots are computed
by Lemma~\ref{C:function-field-pth-roots}.
Each successful step lowers the degree by $p$, so the
procedure stops, and the resulting polynomial is
the unique irreducible factor just proved to exist.
Recover its multiplicity in the original $f$ by exact
division. Repeating for all factors of $f^{p^e}$ gives
the entire factorization of $f$.

This handles a general simple step by first its separable
intermediate field and then its purely inseparable part.
It uses factorization only over $E$ and the separable
extension of $E$ constructed by the preceding step.
Thus it is a genuine induction on the supplied tower,
not an appeal to factorization over the field currently
being constructed. All irreducibility checks needed
to verify additional finite tower presentations are
provided by the same induction. This proves the theorem.
\end{proof}

\begin{theorem}[Proved: complete coefficient-algebra algorithms over
imperfect algebraic function fields]
\label{C:function-field-WG}
Fix an odd prime $p$. Let $K$ be any field supplied as
an explicit finite tower over
$\overline{\mathbf F}_p(t)$, and let $R$ be a
finite-dimensional commutative $K$-algebra supplied
by a multiplication table. For every $m\geq0$,
there are ordinary terminating algorithms giving
canonical indices for all indecomposable finite
projective rational $O_{m,R}$-modules, the explicit
additive form coordinates of
Theorem~\ref{C:artinian_effective_C1}, and equality,
products and every specified exterior coefficient
in the split group $W_G(BO_{m,R})$.
No unclassified isomorphism quotient or unspecified
factorization oracle occurs in this assertion.
The module indices are enumerable natural-number
matrix indices at each finite dimension and coefficient
degree, not a finite list of all dimensions.
\end{theorem}
\begin{proof}
The field $K$ is an infinite countable computable
field, with decidable exact arithmetic and equality
in its tower power basis. It is a function field
of transcendence degree one over the algebraically
closed constant field $\overline{\mathbf F}_p$.
Tsen's theorem therefore says that it is $C_1$
\cite[Theorem 59.67.10, Tag 03RD]{C-StacksTsenFunctionFields}.
This statement does not assume that the displayed
extension over $C(t)$ is separable.

Theorem~\ref{C:function-field-factorization} supplies
factorization over $K$ and every further explicit
finite tower over it, since concatenating the towers
produces another finite tower over $C(t)$.
All hypotheses of
Theorem~\ref{C:artinian_effective_imperfect_C1}
are consequently proved for this actual field family.
Apply its direct matrix construction to $R$.
Projectivity is decided by the finite linear splitting
test; endomorphism radicals and finite residue fields
over $K$ are computed by the terminating tower
certificate; isomorphism and square classes are
decided by the determinant and factorization tests.
The dual is $\operatorname{Hom}_R(-,R)$.
The residual $C_1$ form coordinates, and the balanced
tensor and exterior constructions over $R$, are
exactly those established there. Every occurrence
of its former arithmetic access is now replaced by
the proved algorithms above. This proves all the
claimed outputs.

In particular the family includes the imperfect field
$\overline{\mathbf F}_p(t)$ itself: the valuation
of $t$ at zero is one, while the valuation of a
$p$-th power is divisible by $p$.
The finite extensions remain within the stated
function-field family, and their imperfectness
causes no additional oracle requirement.
This is a termination theorem; no useful complexity
bound is asserted for the finite searches.
\end{proof}
